\documentclass[12pt,a4paper]{article}
\usepackage[utf8]{inputenc}
\usepackage{amsmath,mathrsfs}
\usepackage{amsfonts}
\usepackage{amssymb}
\usepackage{enumitem,graphicx}
\usepackage[left=2cm,right=2cm,top=2cm,bottom=2cm]{geometry}
\author{Léo Brunswic}
\title{On non trivial cycles of the 3$x$+1 problem}
\usepackage{amsthm}
\usepackage{MnSymbol}
\newtheorem{lem}{Lemma}[section]
\newtheorem{cor}[lem]{Corollaire}
\newtheorem{prop}[lem]{Proposition}

\newtheorem{propri}[lem]{Propriété}
\newtheorem{defi}[lem]{Définition}
\newtheorem{theo}{Théorème}

\newtheorem{cexemple}[lem]{Contre-exemple}
\newtheorem{exemple}[lem]{Exemple}
\newtheorem{rem}[lem]{Remarque}
\newtheorem{conjecture}[lem]{Conjecture}
\newtheorem{question}[lem]{Question ouverte}
\newtheorem{notation}[lem]{Notation}
\newtheorem{exercice}[lem]{Exercice}
\newcommand{\fonction}[5]{\displaystyle#1:\begin{array}{l|rcl}
& \displaystyle #2 & \longrightarrow & \displaystyle #3 \\
    & \displaystyle #4 & \longmapsto & \displaystyle #5 \end{array}}
\newcommand{\fonctiondeux}[7]{\displaystyle#1:\begin{array}{l|rcl}
& \displaystyle #2 & \longrightarrow & \displaystyle #3 \\
    & \displaystyle #4 & \longmapsto & \displaystyle #5 \\
     & \displaystyle #6 & \longmapsto & \displaystyle #7\end{array}}
\def\RR{\mathbb{R}}
\def\CC{\mathbb{C}}
\def\NN{\mathbb{N}}
\def\ZZ{\mathbb{Z}}
\def\A{\mathcal{A}}
\def\D{\mathcal{D}}
\begin{document}
\maketitle
\abstract{We present known result on the so-called 3$x$+1 problem also known as Syracuse sequences problem or Collatz problem.  Let $u_0$ be any integer, define the recurring sequence by $u_{n+1}= u_n/2$ if $u_n$ is even and $3u_n+1$ if $u_n$ is odd.  The only known positive cycle is $(1,4,2)$, the conjecture being that the sequence converges to this cycle whatever positive first term is chosen. We present an extension of a method of Terras to show the only cycles are the known one. Leaving open the rest of the conjecture.
}

\section{Known facts and conventions}
A Syracuse is an integer sequence $(u_n)_{n\in\NN}$ such that   $$u_0\in \ZZ, \text{ and }\quad \forall n\in \NN,\quad u_{n+1}=\left\{\begin{array}{ll}
\displaystyle \frac{1}{2}u_{n} & \text{if }u_n \equiv 0~[2]\\
\displaystyle 3u_n+1&\text{if }u_n \equiv 1~[2]\\
\end{array}\right.$$
The sequence is defined {\it a priori} for rational integers but one notices that the definition makes sense in any interal commutative unitary ring $A$ such that $A/2A \simeq \ZZ/2\ZZ$.  Indeed, in such a ring, $2\neq 0$ and if $x\notin 2A$ then $3x+1\in 2A$. Furthermore, if $x\in 2A$, there exists $y\in A$ such that $2y=x$, and two such $y$ are equal since the ring is integral.  In particular, if $A$ is an algebraic integer ring  fully ramified above $2$ or even the 2-adic completion  of $A$, syracuse sequences are well defined. In such a ring, $4\neq 0$, thus $(1,4,2)$ is a cycle.
We notice that if $\bigcap_{n\in \NN} (2A)^n \neq \{0\}$ then there exists an infinite Syracuse sequence in $2A$. We would like to exclude such a situation and thus assume $A$ is noetherian (which thus exclude this situation by Krull intersection Theorem).
In the following, a noetherian ring $A$ such that $A/2A \simeq \ZZ/2\ZZ$ is called a $\ZZ/2\ZZ$-ring. In such a ring, the 2-valuation is well defined $\nu_2(x):= \max \{n\in \NN ~|~ x\in 2^n A\}$, as well as the associated absolute value and the 2-completion of $A$ we will denote by $A_2$. Actually, one can show that for any $\ZZ/2\ZZ$-ring $A$, $A_2$ is naturally isomorphic to the ring of 2-adic integers $\ZZ_2$.

\begin{defi}  Let $A$ be a  $\ZZ/2\ZZ$-ring, define
$$ \fonctiondeux{f_A}{A}{A}{x\in 2A}{3x+1}{x\notin 2A}{x/2} \text{ and }\fonction{T_A}{2A+1}{2A+1}{x}{\frac{3x+1}{2^{\nu_2(3x+1)}}}.$$
\end{defi}
The proof of the following property is left to the reader.
\begin{prop} Let $A$ be a $\ZZ/2\ZZ$-ring, both $f_A$ and $T_A$ are continuous and extends continuously to $A_2$. Furthermore, their extension is $f_{A_2}$ and $T_{A_2}$ respectively.
\end{prop}

\begin{defi} Let $A$ be a $\ZZ/2\ZZ$-ring. 
For $n\in \NN$, define $$\fonction{a_n}{2A+1}{\NN}{x}{\nu_2(3 T^n_A(x)+1)}$$ and 
$$\fonction{p_n}{2A+1}{\NN}{x}{\sum_{k<n}a_k(x)}$$

\end{defi}

\begin{lem}\label{lem:base}Let $A$ be a $\ZZ/2\ZZ$-ring and let $x\in A$, one has : 
$$\forall n\in \NN, \quad 2^{p_n(x)}T_A^n(x)-3^n x = \sum_{k=0}^{n-1} 2^{p_k(x)}\,3^{n-1-k} .$$
\end{lem}
\begin{proof} 
By induction on $n$, for all $x\in 2A+1$, since $T_A(x) = 2^{-a_1(x)}(3x+1)$ and $a_1(x)=p_1(x)$, the right hand side reduces to $1$ and one obtain the very definition of $T_A(x)$.

Assume the formula is true for all $x\in 2A+1$ and some $n \geq 0$, then 
\begin{eqnarray*}
2^{p_{n+1}(x)}T_A^{n+1}(x)-3^{n+1} x &=& 2^{a_{1}(x)} 2^{p_n(T_A(x))} T_A^n( T_A(x))-3^{n+1}x\\&=& 2^{a_1(x)} \left(3^n T_A(x)+ \sum_{k=0}^{n-1} 2^{p_k(T_A(x))}\,3^{n-1-k}  \right) - 3^{n+1}x  \\&=& \sum_{k=1}^{n} 2^{p_{k+1}(x)}3^{n-1-k}+3^n\left( 2^{a_1(x)}T_A(x)-3x\right)
\\&=&\sum_{k=1}^{n} 2^{p_{k+1}(x)}3^{n-1-k}+3^n 
\\&=& \sum_{k=0}^{n} 2^{p_{k+1}(x)}3^{n-1-k}
 \end{eqnarray*}
\end{proof}
\begin{defi}
Let $n\in\NN^*$ and let $a\in (\NN^*)^\ZZ$, define $$ A_{u}^n(a)=\sum_{k=0}
^{n-1}3^k\prod_{i=0}^{n-1-k}2^{a_{u+i}} $$ 
\end{defi}
\begin{lem} Let $A$ be a $\ZZ/2\ZZ$-ring. 
For $n\in \NN$, the maps $$\fonction{RHS_n}{\{p\in \NN^n~|~ p_0=0~\wedge~p_n \nearrow\nearrow \} }{A}{p}{\sum_{k=0}^{n-1}2^{p_k}3^{n-1-k}} $$ 
are injective.
\end{lem}
\begin{theo}\label{theo:rep} The map $$\fonction{\phi}{2A_2+1}{\{p\in \NN^\NN~|~ p_0=0 ~\wedge~ (p_n) \nearrow\nearrow\}}{x}{(p_n(x))_{n\in \NN}} $$
is bijective and its inverse is: 
$$\fonction{\phi^{-1}}{\{p\in \NN^\NN~|~ p_0=0 ~\wedge~ (p_n) \nearrow\nearrow\}}{2A_2+1}{p}{ -\frac{1}{3}\sum_{n = 0}^{\infty}\frac{2^{p_n}}{3^{n}}}.$$
\end{theo}
\begin{proof} Recall that $3\in A_2^{\times}$. From Lemma \ref{lem:base},
 $$\forall x\in 2A+1,\quad  x = -\frac{1}{3}\sum_{k=0}^{n-1}\frac{2^{p_k(x)}}{3^k}+ \frac{2^{p_n(x)}}{3^n}T_A^n(x)$$
and since $\left|\frac{2^{p_n(x)}}{3^n}T_A^n(x)  \right| \leq 2^{-p_n(x)} $ and $p_n(x)\rightarrow +\infty$ as $n\rightarrow +\infty$, we obtain the result. 
\end{proof}
\begin{cor}Let $A$ be a $\ZZ/2\ZZ$-ring and  $x\in A$. The following are equivalent : 
\begin{enumerate}
\item  $\left(a_n(x)\right)_{n\in \NN}$ is periodic of period $n$
\item  $\left(T^n(x)\right)_{n\in \NN}$ is periodic of period $n$
\end{enumerate}
\end{cor}
\begin{proof}
 The only non obvious part is $(i)\Rightarrow (ii)$.  
 From Theorem \ref{theo:rep}, one has 
 $$ x= -\frac{1}{3}\sum_{k=0}^{+\infty} \frac{2^{p_k(x)}}{3^k} = -\frac{1}{3}\sum_{k=0}^{+\infty} \frac{2^{p_k(T^n(x))}}{3^k} = T^n(x)  $$
\end{proof}
\begin{cor}  Let $A$ be a $\ZZ/2\ZZ$-ring and  $x\in A$ such that $\left(T^k(x)\right)_{k\in\NN}$ is $n$-periodic for some $n$.  If $n$ the minimal period then :
$$\forall i,j\in \lsem 1,n\rsem,\quad A_i^n(a(x))-A_j^n(a(x)) \neq 0$$
\end{cor}
\begin{cor}  
\end{cor}
We now concentrate on  $A=\ZZ$ case. 

\begin{lem} Let $x\in \ZZ^*$ such that $\forall n\in\NN, |T^n(x)|\geq |x|$, one has : $$p_n(x) \log(2) - n \log(3)\leq n\log\left(1+ \frac{1}{|x|}\right)  $$
\end{lem}
\begin{proof}
From Lemma \ref{lem:base}, one has 
\begin{eqnarray*}
e^{\log(2)p_n(x)-n\log(3)}T^n(x) - x &=& \sum_{k=0}^{n-1} e^{\log(2)p_k(x)-k\log(3)}
\\
e^{\log(2)p_n(x)-n\log(3)}&\leq &1+\frac{1}{|x|} \sum_{k=0}^{n-1} e^{\log(2)p_k(x)-k\log(3)}
\end{eqnarray*}
Applying Gronwall's Lemma to the sequence $u_k = \exp\left(\log(2)p_k(x)-n\log(3)\right)$ we get 
$$ e^{\log(2)p_n(x)-n\log(3)} \leq \left(1+\frac{1}{x}\right)^n$$
from which the result follows.
\end{proof}
\begin{lem} Let $x\in \NN^*$.

\begin{itemize}
\item If $n\in\NN$ is such that $T^n(x)\leq x$ then 
$$p_n(x)\log(2)-n\log(3) >0$$ 
\item If furthermore $\forall n\in\NN,~T^n(x)\leq x $ then $$\forall n\in \NN,\quad  p_n(x)\log(2)-n\log(3) \geq n \log\left(1+\frac{1}{x} \right).$$ 
\end{itemize}


\end{lem}
\begin{proof}
Proceed the same way as the preceding Lemma.
\end{proof}
\begin{cor} Let $x\in \NN$ such that $T^n(x) \xrightarrow{n\rightarrow +\infty} +\infty$, then $$\limsup_{n\rightarrow+\infty} \frac{p_n}{n}\leq \log_2(3) $$
\end{cor}
\begin{cor}\label{cor:bounds} Let $x\in \NN$ such that $\exists n\in \NN^*,~ T^n(x) = x$, then for such a $n\in \NN^*$, $$\log_2\left(1+\frac{1}{\min_{k\in \NN}(T^k(x))}\right)  \leq \frac{p_n}{n}-\log_2(3)\leq  \log_2\left(1+\frac{1}{\max_{k\in \NN}(T^k(x))}\right) $$
\end{cor}

\begin{theo}[\cite{MR1074572}]
Let $a,b\in \ZZ$ such that $|a\log(2)-b\log(3)|\leq e^{-a/\delta}$ with $\delta \geq 1$ then 
$$a < 2^{800}\log(3)^{25}\delta^{25} $$
\end{theo}

\begin{theo}[Steiner \cite{}] Let $x\in \NN^*$. 
If 
$$\exists n\in\NN^*,\quad T^n(x)=x \text{~and~} \forall k< n,~p_k(x)=k$$
then $x=1$.
\end{theo}

\begin{theo}[Matveev] Let $a,b\in \NN^*$ such that $2^a>3^b$.
 Then :
 $$2^a-3^b \geq \frac{2^a}{e^da^d} $$
 where $d = 3.5\cdot 10^6$.
\end{theo}
\begin{cor} Let $\mathcal P(\Gamma)$ a property defined on cycles $\Gamma$. If
$$\exists \varepsilon>0,\forall \Gamma, \quad \left(\mathcal P(\Gamma)\quad\Rightarrow \quad \mathrm{GCD}(A_i(\Gamma)~:~ i\in \ZZ) \ll  2^{|\Gamma|(1-\varepsilon)}\right) $$ then there exists only finitely many cycles $\Gamma$ such that $\mathcal P(\Gamma)$.

If furthermore, the domination is effective then the cycles can be determined effectively.
\end{cor}
\begin{proof}
 If the property above is satisfyied, then only finitely many couple $(|\Gamma|,n)$ are possible. For each admissible couple $(|\Gamma|,n)$, there are only finitely many $a(\Gamma)$ of type $(|\Gamma|,n)$, and for each such $a(\Gamma)$ there is at most one syracuse cycle of type $a(\Gamma)$. Each step being effective, if the domination is effective then the determination of the $\mathcal P$-syracuse-cycle  is effective.
\end{proof}
\section{Constraints on cycles}
Our starting point is the following.
\begin{prop} Let $x\in \ZZ$ and $n\in \NN^*$ such that $T^n(x)=x$. 
We have : 
$$ 2^{p_n(x)}-3^n = \mathrm{GCD}\left(\sum_{k=0}^{n-1}3^k 2^{p_n(x)-\sum_{i=l-k}^{l} a_i(x) } : l \in 
\lsem 0,n-1 \rsem \right) $$
\end{prop}
\begin{proof}
 From Lemma \ref{lem:base}, 
$$\forall l\in \lsem0,n-1\rsem,\quad\left( 2^{p_n(x)}-3^n\right)T^l(x) = \sum_{k=0}^{n-1}3^k 2^{p_n(x)-\sum_{i=l-k}^{l} a_i(x) }$$
thus 
$$\forall l\in \lsem0,n-1\rsem,\quad\left( 2^{p_n(x)}-3^n\right) ~\left|~ \sum_{k=0}^{n-1}3^k 2^{p_n(x)-\sum_{i=l-k}^{l} a_i(x) }\right.$$
and thus
$$2^{p_n(x)}-3^n ~\left| ~\mathrm{GCD}\left(\sum_{k=0}^{n-1}3^k 2^{p_n(x)-\sum_{i=l-k}^{l} a_i(x) } : l \in 
\lsem 0,n-1 \rsem \right)\right. 
.$$
And since, 
$$3\sum_{k=0}^{n-1}3^k 2^{p_n(x)-\sum_{i=-k}^{0} a_i(x) }-2^{a_1(x)}\sum_{k=0}^{n-1}3^k 2^{p_n(x)-\sum_{i=1-k}^{1} a_i(x) } = 3^n-2^{p_n(x)} $$
we have the result.
\end{proof}

We expect the right hand side to be small and bounds like the one obtained by Baker thus give a lower bound on the GCD side which is roughly of size of $3^n$. So any upper bound on the GCD which is $o(3^n)$ would imply there are finitely many cycles. 

Let us notice there are many "1" among the $a_i(x)$ is $x$ is part of a cycle. 

\begin{prop} Let $x\in \ZZ$  and $n\in \NN^*$ such that $T^n(x)=x$. 
If $|x|>1000$ then : 
$$ \#\left\{i\in \lsem 1,n\rsem ~|~ a_i(x) =1\right\} \geq \alpha n $$
with $\alpha = 2-\log_2(3)-10^{-10}> 0.415037499$

\end{prop}
\begin{proof}
Since the conjecture has been tested for all integers lesser than $10^{10}$ and no cycles have been found containing an integer bigger than $1000$, then $|x|>10^{10}$. Let $a=\#\left\{i\in \lsem 1,n\rsem  ~|~ a_i(x) =1\right\} $, we have $p_n\geq a+2(n-a)$. Therefore, from Corollary \ref{cor:bounds} we  have 
$$ \log(1+1/|x|) \geq \frac{a+2(n-a)}{n} - \log_2(3)$$
and thus :
$$ \frac{a}{n} \geq 2-\log_2(3)-\log(1+1/|x|). $$

\end{proof}
\begin{lem} Let $n,p\in\NN^*$ with $p< n$ and let $a \in (\NN^*)^\ZZ$ $n$-periodic, if there exists $i,j\in \lsem 0,n-1\rsem$, such that $$i\neq j\quad \text{and}\quad \forall k\in\lsem 1,p\rsem,~  a_{i+k}=a_{j+k} $$
then :
$$ \nu_2(A_i-A_j)\geq \sum_{k=1}^p a_k \quad \text{and}\quad  \nu_3(A_{p+i}-A_{p+j})\geq p  $$

$$A_i- A_j \neq 0 $$
\end{lem}
\begin{defi}[$\varepsilon$-circuit] Let $x\in \ZZ$ and $\varepsilon>0$. $\left(x\right)_{n\in\NN}$ is an $\varepsilon$-circuit if the sequence is $n$-periodic for some 
$n\in\NN^*$ and their exists a term $y$ in the sequence such that $$\forall i \in \lsem
 1,p\rsem,~a_i(y)=1$$
for some $p\geq \varepsilon n-1$.
\end{defi}
Notice that a circuit (in the sense Steinitz) is a 1-circuit.
\begin{theo} If $\varepsilon\geq 10^{-14}$, there exists no $\varepsilon$-circuit other than the known ones.
\end{theo}
\begin{proof}
Let $x$ a $\varepsilon$-circuit of minimal period $n\geq1$, let $i\in \NN$ and  let $p\geq \varepsilon n$ such that  
 $$\forall k \in \lsem
 1,p\rsem,~a_{i+k}(x)=1.$$
We consider $A_{p+i}(y)-A_{p+i+1}(y)$.

On the one hand, from the previous Lemma, this difference is divisible by $3^p$, and since $i\neq i+1$ and $n$ is a minimal period we have $A_{p+i}(y)-A_{p+i+1}(y)\neq 0$. 
Furthermore, $2^{p_n}-3^n$ is not divisible by $3$ thus 
$$2^{p_n}-3^n \leq 3^{-p}\left|A_{p+i}(y)-A_{p+i+1}(y)\right|.$$
Furthermore
\begin{eqnarray}
3^{-p}\left(A_{p+i}(y)-A_{p+i+1}(y)\right)&=& \sum_{k=1}^{n-1-p} 3^{n-1-p-k}\left( 2^{\sum_{j=0}^{k-1} a_{p+i}(y)}-2^{\sum_{j=0}^{k-1} a_{p+i+1}(y)} \right) \\
&=&\sum_{k=1}^{n-1-p} 3^{n-1-p-k} 2^{\sum_{j=1}^{k-1} a_{p+i}(y)}\left(2^{a_{p+i}}-2^{a_{p+i+k}(y)} \right) \\
&=&\sum_{k=1}^{n-1-p} 3^{n-1-p-k} 2^{\sum_{j=1}^{k-1} a_{p+i}(y)}\left(2-2^{a_{p+i+k}(y)} \right) \\
&\leq& 3^{n-p}\frac{3^{n\eta}-1}{3^\eta-1}\\
&\leq&  2\log(3)3^{(1-\varepsilon+\eta)n}.
\end{eqnarray} 

On the other hand from Matveev Theorem, one has  
$$\left| 2^{p_n}-3^{n}\right|\geq \frac{2^{p_n}}{(p_ne)^C}$$
with $C = \frac{1}{2}e2^{4.5}30^5\log(3) \simeq 3.6\cdot 10^{8}$.

Since $2^{p_n}\geq 3^n$ and $p_n\leq n\log(3+\eta)$ we thus obtain 
$$  3^{n(\epsilon-\eta)} \leq 2 \left(\log(3+\eta)n\right)^{3.7\cdot 10^8}.$$


\end{proof}



\begin{eqnarray}
3^{-p}\left(A_{p+i}(y)-A_{p+i+1}(y)\right)&=& \sum_{k=1}^{n-1-p} 3^{n-1-p-k}\left( 2^{\sum_{j=0}^{k-1} a_{p+j}(y)}-2^{\sum_{j=0}^{k-1} a_{p+j}(y)} \right) \\
&\leq&\sum_{k=1}^{n-1-p}3^{n-1-p-k}\left(2^{k(\log_2(3+\eta)}-2^{k}\right)\\
&\leq&\sum_{k=1}^{n-1-p}3^{n-1-p}3^{k\eta} \\
&\leq& 3^{n(1-1/\log(n))}\frac{3^{(1-1/\log(n))n\eta}-1}{3^\eta-1}\\
&\leq&  2\log(3)3^{(1-\varepsilon+\eta)n}.
\end{eqnarray} 











\section{Upper bounds}
Let $S := \{2^a3^b ~:~ a,b\in \ZZ\}$
\begin{defi} Let $\Gamma=(x_i~:~ i \in \ZZ)$ be syracuse cycle of period $n>0$, denote $\forall i\in\ZZ , a_i :=T^i(x_0)$.
 We denote by $\mathcal A_\Gamma$ the set of ascending term  in a give cycle $\Gamma$ and by $\mathcal D_\Gamma$ the set of descending terms : 
 $$\A_\Gamma = \{i\in \ZZ~|~ a_i=1\} \quad\quad \D_\Gamma = \{i\in \ZZ~|~ a_i>1\}$$
\end{defi}
\begin{lem}
Let $(\Gamma,a)$ be a Syracuse cycle and let $i\in\ZZ$, then 
$$ \Delta_i := A_{i+1}-A_{i} = \sum_{k=1}^{n-1}3^{n-1-k}2^{a_i+\cdots+a_{i+k-1}} \left(2^{a_{i+k}}-2^{a_i}\right)$$ 

In particular, if $a_i=1$ then :
$$ \Delta_i = \sum_{k=1,k\in\D}^{n-1}3^{n-1-k}2^{a_i+\cdots+a_{i+k-1}} \left(2^{a_{i+k}}-2\right)$$ 
\end{lem}
 
\begin{lem}
Let $(\Gamma,a)$ be a Syracuse cycle, let $m = \mathrm{GCD}(A_i ~:~ i\in \ZZ)=2^{p_n}-3^n$ and let $i\in\ZZ$, then, : 

$$a_i=a_{i+1}=1 \quad\Rightarrow \quad\Delta_{i} \sim_S \Delta_{i+1}$$
\end{lem}

\begin{cor}
If $\A$ and $\D$ are intertwined.
$$ m = A_0\bigwedge_{i=\D}\Delta_{i-1} $$
\end{cor}

\begin{theo}[Steiner] Let $\Gamma$ be a Syracuse n-cycle.
 If $\#\D_\Gamma=1$ then $\Gamma = (\cdots,1,\cdots)$.  
\end{theo}
\begin{proof}
Assume $\D=n\ZZ$, from the preceding corollary,
$m = A_0 \wedge \Delta_{-1}$. We write $\alpha = a_0-1 = \frac{1}{2}(|a(\Gamma)|-n)$
In particular we have~:
$$\Delta_1 \sim_S  2^{\alpha}-1$$
But since $2^\alpha = 2^{|\Gamma|-n}\leq 2^{|\Gamma|(1-\log_3(2))}$ the hypothesis to the corollary to Matveev Theorem are satisfied (effectively) and the finitely many cycle $\Gamma$ such that $\#\D_\Gamma=1$ can be effectively determined. With the aid of a computer, one  can check that the only one is $(\cdots,1,\cdots)$.
\end{proof}


\begin{theo} Let $\Gamma$ be a Syracuse $n$-cycle, then $\#(\D_\Gamma\cap \lsem 1,n\rsem)\neq 2$.
\end{theo}
\begin{proof} Let $k_1,k_2\in\lsem 1,n\rsem$ such that $\#\D_\Gamma=(k_1+n\ZZ)\cup(k_2+n\ZZ)$.  Consider $\Delta_i$ and $\Delta_j$ for $i,j\in \ZZ$ in such a way that $a_i=a_j=1$, $k_2-n< j<k_1<i < k_2$. 
We have : 
\begin{eqnarray}
\Delta_i&=&  3^{n-1-(k_1-i)}2^{k_1-i}\left(2^{a_{k_1}}-2\right)
	+3^{n-1-(k_2-i)}2^{k_2-i+a_{k_1}-1}\left(2^{a_{k_2}}-2\right)\\
\Delta_j&=&  3^{n-1-(k_2-j)}2^{k_2-j}\left(2^{a_{k_2}}-2\right)
+3^{n-1-(k_1+n-j)}2^{k_1+n-j+a_{k_2}-1}\left(2^{a_{k_1}}-2\right)
\end{eqnarray} 
We write $\delta_1 = k_2-k_1 $, $\delta_2 = n+k_1-k_2$, $\alpha_1 = a_{k_1}-1$ et $\alpha_2=a_{k_2}-1$ so that we get : 
\begin{eqnarray}
\Delta_i&\sim_S&  3^{\delta_2}\left(2^{\alpha_1}-1\right)
	+2^{\delta_2+\alpha_1}\left(2^{\alpha_2}-1\right)\\
\Delta_j&\sim_S&  3^{\delta_1}\left(2^{\alpha_2}-1\right)
	+2^{\delta_1+\alpha_2}\left(2^{\alpha_1}-1\right)\\
\end{eqnarray} 
Without loss of generality we may assumer $\delta_1\leq\delta_2$. We thus have $\delta_1\leq n/2$ and $\alpha_1+\alpha_2 = |\Gamma|-n$, and thus : 
$$ \Delta_j \leq 3^{n/2}2^{|\Gamma|-n}+ 2^{|\Gamma|-n}2^{n/2}\leq 2^{|\Gamma|-(\log_3(2)-1/2)|\Gamma|}\leq 2^{|\Gamma|(1-\varepsilon)}$$
with $\varepsilon=\log_3(2)-1/2$.
One can then use the corollary to Matveev Theorem and check the finitely many remaining cases.

It remains to consider the case where $k_2=k_1+1$. 
Consider
\begin{eqnarray}
A_{k_2}-A_i &=& \sum_{k=0}^{n-1}3^{n-1-k}2^{a_i+\cdots a_{i+k_2-1}}
\end{eqnarray}

\end{proof}

\section{Cyclotomic arithmetic}
We now wish to study the GCD problem to obtain generic upper bounds. 
We may notice that a the $A_i$ can be written as convolution product of $(3^k)_{k\in\lsem 1,n\rsem}$ with $\left(2^{p_{k+i}-p_i}\right)_{k\in \lsem1,n\rsem}$. 
To begin with 
\begin{eqnarray}
\left(2^{|\Gamma|}-3^n\right)\begin{pmatrix}
I_0 \\ \vdots \\ I_n 
\end{pmatrix} &=& \begin{pmatrix}
&&\\
&2^{p_{i+j}-p_i}&\\
&&
\end{pmatrix} \begin{pmatrix}
3^{n-1}\\ \vdots \\ 1
\end{pmatrix}
\\
\left(2^{|\Gamma|}-3^n\right)\begin{pmatrix}
 \vdots \\ 2^{p_{i}-i|\Gamma|/n}I_i \\ \vdots
\end{pmatrix}&=&\begin{pmatrix}
&&\\
&2^{p_{i+j}-(i+j)|\Gamma|/n}&\\
&&
\end{pmatrix} 
\begin{pmatrix}
\vdots \\ 2^{i|\Gamma|/n}3^{n-i}\\ \vdots
\end{pmatrix}
\\
\left(2^{|\Gamma|}-3^n\right)\begin{pmatrix}
 \vdots \\ 2^{p_{i}-i|\Gamma|/n}I_i \\ \vdots
\end{pmatrix}&=&\begin{pmatrix}
\vdots\\
2^{p_{i}-i|\Gamma|/n}\\
\vdots
\end{pmatrix} 
\star
\begin{pmatrix}
\vdots \\ 2^{(n-i)|\Gamma|/n}3^{i}\\ \vdots
\end{pmatrix}
\end{eqnarray} 
Therefore, using discrete Fourier transform we have :
$$\left(2^{|\Gamma|}-3^n\right)\mathcal F\begin{pmatrix}
 \vdots \\ 2^{p_{i}-i|\Gamma|/n}I_i \\ \vdots
\end{pmatrix}=\mathcal F\begin{pmatrix}
\vdots\\
2^{p_{i}-i|\Gamma|/n}\\
\vdots
\end{pmatrix} 
\mathcal F
\begin{pmatrix}
\vdots \\ 2^{(n-i)|\Gamma|/n}3^{i}\\ \vdots
\end{pmatrix}
$$ 

The first coordinate of which gives:
 $$\left( 2^{|\Gamma|}-3^n\right)\sum_{k = 0}^{n-1} I_k 2^{p_k-k|\Gamma|/n}  = 
 \sum_{k=0}^{n-1}2^{p_{k}-k|\Gamma|/n}  \sum_{k=0}^{n-1}3^{k}2^{(n-k)|\Gamma|/n}  
 =\frac{3^{n}-2^{|\Gamma|}}{3-2^{|\Gamma|/n}}\sum_{k=0}^{n-1}2^{p_{k}-k|\Gamma|/n} $$
Hence, 
 $$\left( 2^{|\Gamma|/n}-3\right)\sum_{k = 0}^{n-1} I_k 2^{p_k-k|\Gamma|/n}   
 =\sum_{k=0}^{n-1}2^{p_{k}-k|\Gamma|/n} $$
Thus:
$$\sum_{k = 0}^{n-1} 2^{p_k-k|\Gamma|/n}\left(I_k\left( 2^{|\Gamma|/n}-3\right)-1\right)   
 =0$$
 Define $\omega$ a primitive $n$-th root of unity and $b:=2^{1/n}$ and 
 $K=\mathbb Q[\omega,b]$.
 We have :
 $$\sum_{k = 0}^{n-1} b^{np_k-k|\Gamma|}\left(I_k\left( b^{|\Gamma|}-b^n-1\right)-1\right)   
 =0$$
\bibliography{ref}	
\end{document}
